\section*{الفصل \ref{ch:g2:measuring-time} --- قياس الزمن}

\begin{solution}{exo:g2:measuring-time:1}
الاستيقاظ، فالفطور، فلعب بعد الظهر، فالعشاء، فالنوم.
\end{solution}

\begin{solution}{exo:g2:measuring-time:2}
لحظة؛ \omterm{def:g2:measuring-time:duration}{مدة}؛ لحظة؛ \omterm{def:g2:measuring-time:duration}{مدة}.
\end{solution}

\begin{solution}{exo:g2:measuring-time:3}
الشمعة الزرقاء دامت أطول. فقد بدأت الشمعتان في اللحظة نفسها،
والتي بقيت مشتعلة حين انتهت الأخرى هي الأطول دوامًا --- وهذه
قاعدة السباق العادل بعينها.
\end{solution}

\begin{solution}{exo:g2:measuring-time:4}
في كل مرة تُقلب يستغرق الرمل \emph{\omterm{def:g2:measuring-time:duration}{المدة} نفسها} لينسكب من الخصر
الضيّق. والحَكَم الذي لا يتغيّر عادل مع الجميع.
\end{solution}

\begin{solution}{exo:g2:measuring-time:5}
إحساسنا بالزمن ساعة رديئة: فالانتظار المملّ يزحف والمتعة تطير.
أما الساعة فتعدّ بانتظام مهما شعر بن، فالانتظار كان أقصر حقًّا
من الفيلم --- وإنما \emph{بدا} أطول.
\end{solution}

\begin{solution}{exo:g2:measuring-time:6}
الأجوبة مختلفة؛ والمعتاد بضع ثوانٍ للاسم ونحو $20$ ثانية
للحذاء. ويجب أن يبقى العدّ هادئًا منتظمًا.
\end{solution}

\begin{solution}{exo:g2:measuring-time:7}
العطسة: ثوانٍ. وليلة النوم: ساعات. وربط رباط الحذاء: ثوانٍ.
والصباح المدرسي: ساعات.
\end{solution}

\begin{solution}{exo:g2:measuring-time:8}
عند الخامسة تمامًا يشير العقرب القصير إلى $5$ (والعقرب الطويل
إلى الأعلى تمامًا)؛ وعند الخامسة والنصف يشير العقرب القصير إلى
\emph{منتصف الطريق} بين $5$ و$6$ (والعقرب الطويل إلى الأسفل
تمامًا).
\end{solution}

\begin{solution}{exo:g2:measuring-time:9}
أجرِ كلًّا منهما في سباق مع الساعة الرملية نفسها: عند الحمّام
أدِر الساعة وعُدّ كم دورة كاملة منها دام الحمّام؛ وافعل مثل ذلك
غدًا مع الدشّ. ثم اكتب العددين (مثلًا: ``الحمّام: $4$ دورات،
والدشّ: $2$'')؛ فالدورات الأكثر تعني الأطول، ويستطيع كل من عنده
الساعة نفسها أن يتحقّق.
\end{solution}

\begin{solution}{exo:g2:measuring-time:10}
توم أقرب. فعدّ الثواني لا يكون أمينًا إلا بإيقاع هادئ منتظم ---
نحو نبضة قلب هادئة لكل عدد، أو ``تمساح صغير'' لكل عدّة. أما
الإسراع في الأعداد فيجعل الثواني أقصر مما ينبغي.
\end{solution}
